\documentclass[journal,10pt]{IEEEtran}
\usepackage{amsmath,amssymb}
\usepackage{booktabs}
\usepackage{array}
\usepackage{graphicx}
\usepackage[dvipsnames]{xcolor}
\usepackage[colorlinks=true,linkcolor=NavyBlue,citecolor=PineGreen,urlcolor=NavyBlue]{hyperref}
\usepackage{xspace}
%% Generated by scripts/make_agentic_artifacts.py; do not edit.
%% Agentic Communication prose numbers must use these macros rather than literals.
\newcommand{\agenticDiagnosisCapabilityRequests}{4.00}
\newcommand{\agenticDiagnosisContexts}{2.00}
\newcommand{\agenticDiagnosisMaterializedBytes}{65046.8}
\newcommand{\agenticDiagnosisModelRequests}{2.00}
\newcommand{\agenticDiagnosisPercepts}{4.00}
\newcommand{\agenticEvidenceAwareExtraCapabilities}{0.00}
\newcommand{\agenticFixedOrderCapabilityRequests}{28.00}
\newcommand{\agenticFixedOrderMaterializedBytes}{480537.8}
\newcommand{\agenticFixedOrderModelRequests}{14.00}
\newcommand{\agenticFullDumpProtocolBytes}{31293.4}
\newcommand{\agenticGenericReactCapabilityRequests}{0.00}
\newcommand{\agenticGenericReactContexts}{0.00}
\newcommand{\agenticGenericReactMaterializedBytes}{-377527.4}
\newcommand{\agenticGenericReactModelRequests}{0.00}
\newcommand{\agenticGenericReactProtocolBytes}{26943.2}
\newcommand{\agenticGlobalBackupBytes}{2141.6}
\newcommand{\agenticGlobalCollection}{0.51062}
\newcommand{\agenticGlobalEnergy}{0.35077}
\newcommand{\agenticGlobalLatencyP}{468.0}
\newcommand{\agenticGlobalTdr}{0.32436}
\newcommand{\agenticLocalCollection}{0.61124}
\newcommand{\agenticLocalConfirmRate}{85.00}
\newcommand{\agenticLocalDumpBytes}{29575.6}
\newcommand{\agenticLocalDumpEvidence}{27.32}
\newcommand{\agenticLocalLatencyP}{3000.0}
\newcommand{\agenticLocalSufficiency}{1.000}
\newcommand{\agenticLocalTaskBytes}{19763.6}
\newcommand{\agenticLocalTaskEvidence}{10.32}
\newcommand{\agenticLocalTdr}{0.45020}
\newcommand{\agenticRobustnessAxes}{5}
\newcommand{\agenticRobustnessCoordinates}{10}
\newcommand{\agenticSourceHarvestTwentyFour}{0.09480}
\newcommand{\agenticSourceHarvestTwentyThree}{0.14674}
\newcommand{\agenticSourceHarvestTwentyTwo}{0.14604}
\newcommand{\agenticTaskContextProtocolBytes}{31293.4}

% Generated by scripts/make_agentic_paper_v1.py; do not edit.
\newcommand{\AgenticMainMethodEpisodes}{15}
\newcommand{\AgenticMainPhysicalExactEpisodes}{15}
\newcommand{\AgenticWOATokenReductionOverall}{35.804\%}
\newcommand{\AgenticQPDeepSeekPhysical}{5/5 execution-equivalent (4 direct + 1 recovered)}
\newcommand{\AgenticQPMiMoPhysical}{5/5 direct}
\newcommand{\AgenticQPQueries}{16}
\newcommand{\AgenticQPGain}{TDR +5.238 pp; AoI -1150.7 s}
\newcommand{\AgenticWOATaskReductions}{48.9\%, 44.6\%, 20.9\%}

\raggedbottom

\title{\textbf{Decision-Semantic Context for Agentic Communication\\
under Intermittent Connectivity}}

\author{Working draft --- \texttt{adaptive-communication} repository}
\date{}

\begin{document}
\maketitle

\begin{abstract}
LLM agents for communication systems are often evaluated on a telemetry snapshot: the model sees
network state, chooses a tool, and receives the result. Intermittent disaster-monitoring networks
violate this abstraction. Evidence is owned by different network locations, acquiring it uses the
same impaired communication substrate, candidate actions become live or retire over time, and an
accepted action may remain in flight after the model call ends. We introduce a
\emph{decision-semantic compiler} that maps Operational Task, Evidence, Capability/Authority, and
execution state into a dynamic plan--evidence relation. The runtime exposes only control-live
dependencies to the model, represents action and no-action sufficiency explicitly, acquires owner
evidence when a live plan remains blocked, and persists the resulting semantic commitment through
asynchronous physical execution. In five-seed query-negative experiments, the Method preserves the
deterministic physical reference in \AgenticMainPhysicalExactEpisodes/\AgenticMainMethodEpisodes{}
episodes. Against a same-interface WirelessOpsAgent-style assurance adaptation, reliability ties
while the Method uses \AgenticWOATokenReductionOverall{} fewer model tokens overall. A
query-positive gateway-backup family closes the full loop---blocking EvidenceNeed, owner query,
guard closure, commitment, and physical effect---with both DeepSeek Flash and MiMo v2.6 Flash,
yielding \AgenticQPGain{} versus acquisition disabled. We delimit the contribution against
decision-region determination, action-sufficient representation, stochastic Boolean evaluation,
provenance witness selection, and semantic communication: the claim is a dynamic
compiler/runtime semantics for communication-constrained evidence and execution, not a new
generic theory of active information acquisition.
\end{abstract}

\begin{IEEEkeywords}
Agentic communication, disaster monitoring, LLM agents, intermittent connectivity, evidence,
context, LoRaWAN, resilient sensing, benchmark.
\end{IEEEkeywords}

\section{Introduction}
Future wireless systems are moving from task-specific learning toward foundation models and
agentic autonomy~\cite{liang2025llmwireless,sheng2025wirelessfm,jiang2026largemodelsurvey,lu2026agenticgnn}.
Recent work already explores intent-driven network orchestration, agentic physical-layer control,
LLM-assisted network optimization, and disaster-oriented UAV coordination~\cite{jiang2026agenticibn,li2026agenticwireless,topollm,xu2025uavllm,sharma2024uavmarl}.
These systems largely assume that the state needed for a decision is available to the agent when
reasoning begins. That assumption is fragile in long-lived field monitoring. A useful fact may
reside at a sensor or gateway; the path required to retrieve it may be delayed or unavailable; and
the same scarce communication opportunity may also be needed to install the corrective action.

We study this coupling in pre-disaster geohazard monitoring. The inherited system model contains
fourteen battery/solar sensor nodes, LoRaWAN Class~A access, a field gateway, cellular primary
backhaul, sparse backup communication, node caches, and remote sampling/report configuration. It
also keeps an evidence discipline already established in the systems version of the repository:
an access receipt is not backhaul reachability, and backhaul reachability is not central
completion. The Agent does not predict a landslide or invent a warning level. External authority
provides an Operational Task; the Agent must execute that task while legal evidence and physical
actions are constrained by the communication system itself.

This setting connects several mature literatures. Delay-tolerant networking treats disconnection,
store-carry-forward, finite buffers, and time-varying paths as first-class system properties
~\cite{fall2003dtn,jain2004dtnrouting,dtnarch,rfc5050,rfc9171,iranmaneshdtnqueue}. Age of
Information and wireless networked control emphasize that freshness, delay, control utility, and
communication scheduling are coupled~\cite{kaul2012aoi,yates2021aoisurvey,bacinoglu1905,wang2021voi,park2018wncs}.
Semantic communication further argues that information value is task dependent rather than equal
to raw data volume~\cite{uysal2021semantic,guo2025semcomnet,jiang2025worldmodelsemcom}.
Our question is narrower: how should an LLM-facing runtime maintain the current relation between
plans and evidence while those plans, observations, and physical effects evolve?

The paper is organised around two interfaces and one compiler:
\begin{equation}
\boxed{\text{Physical/Data Plane}\rightarrow\text{Evidence World}}
\label{eq:method-evidence}
\end{equation}
\begin{equation}
\boxed{\text{Operational Task}\rightarrow\text{Runtime TaskContract}}
\label{eq:method-runtime}
\end{equation}
and a decision-semantic compiler between them. The compiler constructs candidate plans, traces
plan-local dependencies, retires relations that are no longer control-live, projects the
model-facing decision surface, and determines whether the current state is sufficient for action,
sufficient for no action, or blocked on an owner-scoped EvidenceNeed.

\subsection{Contributions and claim boundary}
The current paper makes four concrete contributions.
\begin{enumerate}
\item \textbf{Decision-semantic compilation.} Typed Task, Evidence, Capability/Authority, and
execution contracts are compiled into candidate plans, plan-local dependencies, live/retired
relations, and explicit action/no-action sufficiency instead of being passed to the model as an
undifferentiated telemetry dump.
\item \textbf{Closed-loop evidence acquisition and commitment.} A blocked live plan exposes a
legal owner query; returned evidence can reject or support the plan, after which the semantic
commitment persists through requested/accepted/delivered/applied/confirmed physical execution.
\item \textbf{Communication $\times$ Agent evaluation with strong boundaries.} R0--R3 replay,
model-call ledgers, deterministic compilers, ordinary communication controls, evaluator-only
oracles, and a same-interface WirelessOpsAgent-style assurance adaptation separate model behavior
from physical consequence.
\item \textbf{Source-grounded reproducibility.} O1--O6 tasks, source-period splits, held-out
Task/source transfer, mechanism interventions, and all main paper tables are generated from frozen
contracts/results; the claim ledger and compact paper result are the numeric authorities.
\end{enumerate}
We do not claim a new generic principle of decision-aware acquisition, minimal sufficient state,
minimum certificate selection, provenance witness minimization, or value of information. Those
problems are covered by established work~\cite{javdani2014drd,huang2022asr,deshpande2013sbfe,hellerstein2022adaptivity,green2007provenance,lin2026blip}.
The frozen real-model claim ceiling is A7--A11 in the repository claim ledger.

\section{Scenario and System Model}
\label{sec:system}
We consider pre-disaster mountain monitoring with nodes $i\in\mathcal N$, a gateway $g$, and a
centre $c$. Time is slotted. A monitoring obligation
\begin{equation}
o=(i,m,r_o,[s_o,e_o],d_o)
\end{equation}
requires measurand $m$ from node $i$ to be sampled within $[s_o,e_o]$ and delivered centrally by
$d_o$. Its success indicator is $I_o^{\rm succ}$ and the primary service metric is
\begin{equation}
\mathrm{TDR}=\frac{1}{|\mathcal O|}\sum_{o\in\mathcal O} I_o^{\rm succ}.
\label{eq:tdr}
\end{equation}
Collection, missing delivery, censoring, latency, Age of Information (AoI), recovery and resource
metrics are reported separately; the Agent never removes an obligation from the denominator.

Each node maintains battery $B_{i,t}$, source cache $Q_{i,t}$ and installed configuration
$C_{i,t}$. Battery evolution follows
\begin{equation}
\begin{aligned}
B_{i,t+1}=\min\!\Bigl\{B_i^{\max},\bigl[&B_{i,t}+H_{i,t}
-E^{\rm sense}_{i,t}-E^{\rm UL}_{i,t}\\
&-E^{\rm DL}_{i,t}-E^{\rm DtS}_{i,t}-E^{\rm other}_{i,t}\bigr]^+\Bigr\},
\end{aligned}
\label{eq:battery}
\end{equation}
where $H_{i,t}$ can be driven by source-period-separated NASA POWER irradiance. The deployment
contract follows the monitoring/communication boundary inherited from the systems manuscript and
public geohazard/BeiDou specifications~\cite{dzt0450,dzt0460,bdsps3,ccsdssabr}. A successful
Class~A uplink opens bounded receive opportunities; a centre-originated configuration therefore
requires both a return/control path and a valid access opportunity. Gateway backup, terminal DtS,
and access-assist reuse their existing physical implementations and counters rather than Agent-
reported success.

The complete physical state is $x_t$ and exogenous process is $\xi_t$:
\begin{equation}
x_{t+1}=F(x_t,a_t,\xi_t;\phi).
\label{eq:physics}
\end{equation}
The Agent does not observe $x_t$. At owner/location $\ell$, only a legal observation
$z_t=H(x_{\le t},a_{<t},\xi_{\le t},\ell)$ can enter the runtime.

\section{Data and Evidence World}
\label{sec:evidence}
A typed evidence item is
\begin{equation}
e_j=(p_j,v_j,s_j,o_j,t_j^{\rm gen},t_j^{\rm obs},\rho_j,\nu_j,\sigma_j),
\label{eq:evidence}
\end{equation}
containing proposition/value, source/provenance, owner, generation/observation times, revision,
freshness semantics and status. Status distinguishes at least CURRENT, STALE, NONE\_RECENT,
UNREACHABLE, TIMEOUT and NEGATIVE observation. These distinctions prevent an absent packet from
being silently interpreted as a negative fact.

The Evidence World at time $t$ is
\begin{equation}
\mathcal E_t=\{e_j:t_j^{\rm obs}\le t,\ e_j\text{ passes the runtime gate}\}.
\end{equation}
Evidence revision changes only when factual content changes; age and AoI are derived at
materialisation time. Owner-scoped gateway reads cross the same backhaul abstraction: when that
path is unavailable the Agent receives UNREACHABLE rather than gateway truth.

The primary benchmark uses a source-grounded hybrid simulator: SRTM/ITM terrain propagation,
NASA POWER irradiance/temperature, ChirpBox-derived temporal link processes, public communication
and geohazard-monitoring constraints, and the frozen event-driven monitoring simulator. Source
years are split before random seeds; the current manifest uses 2022/2023/2024 as train/dev/test
source periods. Hazard physics may later generate Operational Tasks, but landslide prediction is
outside the Agent control loop.

\section{Agent Abstraction and Runtime}
\label{sec:runtime}
We separate two meanings of ``task.'' An \emph{Operational Task} is the benchmark/business
objective: target nodes/measurands, monitoring profile, windows/deadlines, authorised effects and
scoring profile. A \emph{Runtime TaskContract} is an Agent program object with desired state,
evidence contract, temporal contract, effect ceiling and completion predicate.

For a runtime task $T_j$, the ContextManifest is
\begin{equation}
C_t=M(T_j,S_t,\mathcal E_t,K_t),
\end{equation}
where $S_t$ is InvestigationState and $K_t$ is the visible capability catalogue. The model input
is $P_t=\mathrm{Materialize}(C_t)$; a planner emits
\begin{equation}
(q_t,a_t,\Delta S_t,stop_t)=\pi(T_j,P_t,K_t).
\label{eq:planner}
\end{equation}
Capabilities split into evidence-use $\mathcal K^E$ and device-action $\mathcal K^A$ but share one
contract surface. An action lifecycle distinguishes requested, accepted/refused, delivered,
applied and confirmed. Hence lack of confirmation never implies that an action was not applied.

The live runtime supports multi-round execution:
\begin{equation}
C_t^{(k)}\rightarrow q_t^{(k)}\rightarrow \mathcal E_t^{(k+1)}
\rightarrow C_t^{(k+1)}\rightarrow a_t.
\label{eq:multiround}
\end{equation}
In particular, a planner can request gateway-owner evidence, trigger an Evidence World revision,
then decide whether to configure a node or invoke a communication fallback.


\subsection{Decision-semantic compiler}
Let $T_t$ be the current Runtime TaskContract, $E_t$ the legal Evidence World, $K_t$ the
capability/authority surface, and $X_t$ the persistent execution state. The compiler constructs a
finite candidate set
\begin{equation}
\mathcal P_t=G(T_t,E_t,K_t,X_t)
\end{equation}
and a typed plan--evidence relation $R_t$. For a plan $p$, the relation records its effect and
scope, authority, feasibility, support references, unresolved guards, and the EvidenceNeeds that
can change its current feasibility. Plans are classified as supported, conditional, rejected, or
dominated.

We distinguish three projections of the same relation. $R_t^{\rm audit}$ retains the complete
candidate/dependency history for replay. $R_t^{\rm control}$ keeps candidates that can still
change the current control decision. $R_t^{\rm model}$ is the materialized control-live surface
shown to the model. Consequently, an unresolved proposition does not remain model-visible merely
because it is unknown: when its plan is rejected or dominated, that dependency retires from the
current decision surface while remaining auditable.

Decision sufficiency is evaluated on $R_t^{\rm control}$. The runtime exposes four operational
states:
\begin{equation}
\begin{aligned}
\mathcal S_t\in\{&\text{sufficient-action},\text{sufficient-no-action},\\
&\text{blocked},\text{undetermined}\}.
\end{aligned}
\end{equation}
A unique supported effect plan can be committed when no EvidenceNeed blocks it. A zero-effect
hold is sufficient only when no other live control-eligible plan can still become a required
effect. Otherwise, an open EvidenceNeed names the legal owner/capability that may resolve the
blocking guard. Returned evidence is recompiled rather than appended as unstructured history.

A semantic commitment is not a one-tick answer. Once selected, the plan is expanded into typed
capability invocations and may progress through requested, accepted, delivered, applied, and
confirmed stages. The planner is re-entered when the commitment basis changes---for example, a
Task revision, guard resolution, authority change, target satisfaction, or execution failure---not
simply because telemetry was refreshed. This separation is what couples Context management to
asynchronous communication execution.

\section{Benchmark Tasks and Experimental Protocol}
\label{sec:protocol}
The benchmark contains six Operational Task families: O1 monitoring continuity, O2 risk
escalation, O3 backhaul-outage sustainment, O4 energy-constrained monitoring, O5 recovery and
reconciliation, and O6 compound long-horizon operation. The same deployment, physical transition,
scorer and information boundary are reused across methods.

We use four replay levels. R0 replays protocol/state transitions; R1 substitutes a planner on
frozen model inputs; R2 reconstructs exact Context materialisation; R3 reruns the full simulator.
The model-facing protocol records ModelRequest, ModelAttempt, ModelUsage and PlannerDecision. R1
diagnostics include stopping correctness, tool selection precision/recall, order LCS and argument
grounding. R3 reports their physical communication consequence.

Online baselines and evaluator-only oracles are typed separately. Current online references
include local/ordinary communication control, deterministic task comply, task-conditioned
evidence, diagnosis-first, fixed-order eager probing, FullDump, and generic-ReAct context.
Delivery/energy oracles are evaluator-only and never appear as online Agent arms. A generic VoI
baseline is intentionally deferred until evidence-acquisition cost/value has a source-backed or
simulator-authoritative scale; arbitrary weights would manufacture a comparison rather than
strengthen it.

Gold replacement follows the cumulative order
\begin{equation}
\begin{aligned}
\text{Task}&\rightarrow\text{EvidenceNeed}\rightarrow\text{Percept}\rightarrow\text{Context}\\
&\rightarrow\text{Selection}\rightarrow\text{Order}\rightarrow\text{Arguments}\\
&\rightarrow\text{Policy}\rightarrow\text{Physical Oracle}.
\end{aligned}
\end{equation}
Runtime-owned upstream replacements require a model rerun; planner-owned selection/order/argument
layers are frozen-input diagnostics; policy/oracle replacements return to R3.

\section{Evaluation}
\label{sec:evaluation}
We report four frozen result groups generated from the compact
\texttt{paper-v1} result authority. Raw live-model traces are retained for local
replay but are not the numeric authority. All model arms use frozen protocol/settings and never
fall back to scripted consumers.

\subsection{Query-negative main table}
Table~\ref{tab:paper-main} compares the Method with task-conditioned, FullDump, and generic-ReAct
inputs on localized O2, O5, and O6. The Method reproduces the deterministic physical reference in
\AgenticMainPhysicalExactEpisodes/\AgenticMainMethodEpisodes{} episodes and commits no extra
observations. This result is deliberately scoped: every Method request in A7 has one ready
supported plan and no visible EvidenceNeed, so the table evaluates compiled control semantics and
execution fidelity rather than autonomous evidence acquisition.
\begin{table*}[t]
\centering\small
\caption{Five-seed query-negative live-model main table. Effect exact is measured against the
frozen candidate/reference semantics; physical exact compares the final simulator trajectory.}
\label{tab:paper-main}
\resizebox{\textwidth}{!}{\begin{tabular}{llrrrrrr}
\toprule
Task & Arm & Effect exact & Err. ep. & Phys. & Obs/ep & Calls/ep & Tok./ep \\
\midrule
Localized O2 & Method & 100.0\% & 0/5 & 5/5 & 0.0 & 6.8 & 36261 \\
Localized O2 & Task-conditioned & 48.8\% & 5/5 & 1/5 & 8.8 & 8.4 & 68107 \\
Localized O2 & FullDump & 50.0\% & 5/5 & 0/5 & 9.2 & 8.0 & 84621 \\
Localized O2 & Generic ReAct & 38.3\% & 5/5 & 0/5 & 9.4 & 9.8 & 114521 \\
O5 & Method & 100.0\% & 0/5 & 5/5 & 0.0 & 9.0 & 61200 \\
O5 & Task-conditioned & 37.3\% & 5/5 & 0/5 & 28.0 & 10.4 & 154886 \\
O5 & FullDump & 44.0\% & 5/5 & 0/5 & 22.6 & 10.2 & 145226 \\
O5 & Generic ReAct & 29.4\% & 5/5 & 0/5 & 17.8 & 11.4 & 131945 \\
O6 & Method & 100.0\% & 0/5 & 5/5 & 0.0 & 8.8 & 100761 \\
O6 & Task-conditioned & 34.9\% & 5/5 & 0/5 & 26.2 & 12.6 & 185166 \\
O6 & FullDump & 27.9\% & 5/5 & 0/5 & 27.6 & 14.0 & 205055 \\
O6 & Generic ReAct & 23.9\% & 5/5 & 0/5 & 7.6 & 9.8 & 113895 \\
\bottomrule
\end{tabular}
}
\end{table*}

\subsection{Same-interface action assurance}
A weak FullDump baseline would not establish whether savings come merely from giving the model
less text. We therefore adapt the assurance pipeline of WirelessOpsAgent~\cite{lu2026wirelessops}
to the same candidate menu, plan-ID expander, capability surface, persistent executor, and model
budget. Table~\ref{tab:paper-woa} shows a reliability tie: the WOA-style proposal is already raw
exact and its repair layer is not activated on these coordinates. The remaining difference is
model-processing cost; the Method reduces total tokens by
\AgenticWOATokenReductionOverall{} across the fifteen episodes. This is a same-interface method
adaptation, not an official reproduction of the authors' implementation.
\begin{table}[t]
\centering\small
\caption{Same-interface WirelessOpsAgent-style comparison.}
\label{tab:paper-woa}
\resizebox{\columnwidth}{!}{\begin{tabular}{lrrrrrr}
\toprule
Task & Method phys. & WOA phys. & WOA raw & Repairs & Method tok. & Reduction \\
\midrule
Localized O2 & 5/5 & 5/5 & 34/34 & 0 & 36261 & 48.9\% \\
O5 & 5/5 & 5/5 & 45/45 & 0 & 61200 & 44.6\% \\
O6 & 5/5 & 5/5 & 44/44 & 0 & 100761 & 20.9\% \\
\bottomrule
\end{tabular}
}
\end{table}

\subsection{Mechanism interventions}
Table~\ref{tab:paper-ablation} separates four mechanisms. FullDump with the same control surface
shows that evidence projection mainly reduces model cost once the decision is already closed.
Removing the explicit no-action sufficiency certificate causes unnecessary investigation in a
closed hold state. The held-out MiMo failure localizes another mechanism: unresolved conditions
belonging to rejected/dominated plans can still distract the model; pruning only those retired
conditions restores exact behavior, while more aggressive pruning provides no extra gain. Finally,
disabling acquisition in the query-positive family prevents the gateway-backup effect and loses
the associated physical improvement.
\begin{table}[t]
\centering\small
\caption{Mechanism interventions; each row changes the stated interface component while keeping
its paired runtime/physical coordinate fixed.}
\label{tab:paper-ablation}
\resizebox{\columnwidth}{!}{\begin{tabular}{p{0.18\columnwidth}p{0.30\columnwidth}p{0.42\columnwidth}}
\toprule
Mechanism & Intervention & Result \\
\midrule
Evidence projection & CF: same candidate, needs, sufficiency + FullDump evidence & O5 action exact=True; O6 hold exact=True; FullDump input increases model cost (see T2) \\
Explicit no-action sufficiency & CS: remove explicit sufficiency from same compact input & O6 hold stop\_exact=False; spurious observations=2 \\
Retired-dependency projection & M keep retired unresolved vs P prune retired-plan unresolved vs D prune more & M effect exact 1/3, obs=4; P 3/3, obs=0; D 3/3, obs=0 \\
Decision-conditioned acquisition & Full loop vs no-acquisition control & 5/5 TDR improved; 5/5 AoI improved; mean TDR +5.238 pp; mean AoI -1150.7 s \\
\bottomrule
\end{tabular}
}
\end{table}

\subsection{Held-out transfer and query-positive acquisition}
Table~\ref{tab:paper-transfer} contains the two tests not covered by the query-negative main
suite. First, a Qili source-derived Operational Task and NASA POWER 2024 test window are held out
from the v6 development table. DeepSeek remains semantically exact; MiMo has one single-turn wobble
but all five episodes preserve final physical fidelity. Second, the gateway-backup family begins
with backup disabled and a delayed primary backhaul. The live plan is conditional until gateway
owner evidence resolves whether received data are actually accumulating behind the primary path.
Both DeepSeek and MiMo perform the legal owner queries, close the guard, commit one backup effect
per seed, and reproduce the deterministic query-positive trajectory. Relative to acquisition
disabled, the mean communication change is \AgenticQPGain{}. DeepSeek seed~0 is the one explicit
recovery case: the episode completed but the original post-processing process was killed after
writing a large trace; streaming recovery reconstructs all 785 assemblies and the query/effect
trace exactly, so the table labels it execution-equivalent rather than pretending that the
original direct physical-equality boolean was persisted.
\begin{table*}[t]
\centering\small
\caption{Held-out task/source/model transfer and the query-positive gateway-backup family.}
\label{tab:paper-transfer}
\resizebox{\textwidth}{!}{\begin{tabular}{p{0.18\textwidth}p{0.12\textwidth}rrrrp{0.20\textwidth}}
\toprule
Setting & Model & Effect exact & Physical & Queries & Gain \\
\midrule
Held-out Qili/NASA-POWER-2024 & DeepSeek Flash & 29/29 & 5/5 & 0 & N/A (query-negative transfer) \\
Held-out Qili/NASA-POWER-2024 & MiMo v2.6 Flash & 29/30 & 5/5 & 1 & N/A (query-negative transfer) \\
Query-positive gateway backup & DeepSeek Flash & 56/56 & 5/5 execution-equivalent (4 direct + 1 recovered) & 16 & TDR +5.238 pp; AoI -1150.7 s \\
Query-positive gateway backup & MiMo v2.6 Flash & 56/56 & 5/5 direct & 16 & TDR +5.238 pp; AoI -1150.7 s \\
\bottomrule
\end{tabular}
}
\end{table*}

\subsection{What the results do and do not establish}
A7 establishes reliable consumption of compiled semantics on query-negative tasks. A8 establishes
a same-reliability model-cost difference against a strong assurance-style interface. A9 provides
a held-out witness for dependency liveness and model transfer. A10/A11 establish that one real
blocking EvidenceNeed can drive an owner query, feasibility transition, semantic commitment, and
physical communication consequence across two model families. None of these results establishes
globally optimal acquisition, an optimal query ordering, or a universal minimum Context.

\section{Related Work}
\label{sec:related}
\subsection{Intermittent, delay-tolerant, and emergency communication}
Our physical setting inherits the challenged-network viewpoint rather than treating outages as
exceptional failures. DTN architecture and routing explicitly model store-carry-forward,
partitioned paths, finite buffers, and delayed delivery~\cite{fall2003dtn,jain2004dtnrouting,dtnarch,rfc5050,rfc9171,iranmaneshdtnqueue}.
Satellite and space-delay-tolerant networking similarly motivate scheduled and non-terrestrial
return opportunities~\cite{ccsdssabr,sateriot}. Mountain monitoring further combines energy,
deadline, and cache constraints studied in energy-harvesting and deadline-aware communication
~\cite{zakeri2311,costa2014pm,kam2016deadline,krishnan2019multihop,leboudec01,wada2009timeout}.
The concrete deployment boundary follows public geohazard-monitoring and BeiDou specifications
~\cite{dzt0450,dzt0460,bdsps3}; the Agent layer does not replace those physical/operational
contracts.

Emergency-network research often restores coverage using UAVs, learned placement, or dynamic
multi-hop networking. Recent work includes MARL-based post-disaster UAV deployment and LLM-guided
multi-hop UAV coordination~\cite{sharma2024uavmarl,xu2025uavllm}, while broader emergency-network
optimization and graph restoration work motivates resilient topology control
~\cite{icldc,icgrestore}. These approaches optimize the communication network itself. Our focus is
different: an LLM-facing runtime must reason over evidence whose availability is produced by that
network and whose actions subsequently change the same physical substrate.

\subsection{Freshness, control value, and semantic communication}
Age of Information formalizes that status utility depends on update timeliness rather than
throughput alone~\cite{kaul2012aoi,yates2021aoisurvey,bacinoglu1905}. Wireless networked control
makes the coupling explicit: communication scheduling, delay, control performance, and lifetime
must be designed jointly~\cite{park2018wncs,wang2021voi}. Existing deadline, retransmission, and
feedback work likewise separates delivery timing from mere packet admission
~\cite{razavivideo,kam2016deadline,wada2009timeout}.

Semantic communication generalizes this observation from freshness to task significance. Uysal
\emph{et al.} argue for joint information generation, transmission, and use under real-time
constraints~\cite{uysal2021semantic}; surveys now extend the paradigm to multi-agent semantic
communication networks~\cite{guo2025semcomnet}. Foundation-model and generative approaches
further explore multimodal knowledge, task-adaptive representation, and world-model-assisted
transmission~\cite{du2024distributedfm,sheng2025wirelessfm,jiang2025worldmodelsemcom}. We borrow
the discipline that information value is task dependent, but do not claim a new VoI or
freshness metric. Our runtime instead gives typed owner/path/freshness semantics to evidence and
asks whether that evidence remains live for the current plan relation.

\subsection{LLMs, foundation models, and agentic wireless systems}
Wireless AI is rapidly moving from task-specific predictors to foundation models and agentic
orchestration. Recent surveys and position papers cover large AI models across PHY, resource
allocation, network management, semantic communication, and agentic systems
~\cite{jiang2026largemodelsurvey,liang2025llmwireless,lu2026agenticgnn}. Shi Jin and collaborators
also demonstrate multi-task wireless foundation models and outline the progression from model
adaptation to autonomous wireless agents~\cite{sheng2025wirelessfm,liang2025llmwireless}.
Agentic physical-layer and intent-based networking frameworks translate high-level goals into
link or slice decisions~\cite{li2026agenticwireless,jiang2026agenticibn}; LLM-based network
planning and restoration systems similarly use models as high-level orchestrators
~\cite{topollm,ossgpt}. DORA provides a complementary disaster-agent benchmark with heterogeneous
data, typed tools, and trajectory diagnostics~\cite{dora2026}.

These systems motivate our model interface, but most begin from an already available network state
or tool result. WirelessOpsAgent moves closer to our concern by building decision-linked evidence,
integrity validation, dependency-scoped repair, revalidation, and an action governor
~\cite{lu2026wirelessops}. We therefore use a same-interface WirelessOpsAgent-style assurance
adaptation as a strong comparison rather than treating evidence validation as our novelty. The
remaining distinction is dynamic: candidate/dependency liveness, separation of audit and
model-facing relations, legal acquisition through the communication path, and persistence of a
semantic commitment while physical execution is still in flight.

\subsection{Decision sufficiency and evidence acquisition}
Decision Region Determination shows that active testing need not identify the full hidden state;
it may stop once all remaining hypotheses lie in an acceptable decision region
~\cite{javdani2014drd}. Action-sufficient state representation similarly studies compact state
information sufficient for downstream control~\cite{huang2022asr}. Stochastic Boolean Function
Evaluation and its adaptivity-gap literature formalize cost-aware testing until a Boolean
certificate is known~\cite{deshpande2013sbfe,hellerstein2022adaptivity}. Database provenance uses
AND/OR derivations and supporting witnesses~\cite{green2007provenance}, and recent verifiable
provenance work directly studies minimal input subsets that preserve LLM-powered data-processing
answers~\cite{lin2026blip}.

These results sharply constrain our claim. We do not introduce decision-aware acquisition,
minimal action-sufficient state, or minimum proof selection. Indeed, a frozen-input audit of the
formal paper workloads found no first-class alternative proof sets, multi-source requirements,
shared blocking EvidenceNeeds, or duplicate evidence keys; a proposed minimum-basis optimizer was
therefore killed before implementation. The contribution retained here is the typed dynamic
compiler/runtime semantics that determines which dependencies are live, when an owner query is
legal, and how the resulting commitment survives asynchronous communication execution.

\section{Limitations}
The current result has five explicit boundaries. First, the query-negative A7 suite contains a
unique ready supported plan on every Method request; it evaluates reliable consumption of compiled
semantics, not difficult multi-candidate search. Second, A10/A11 validate one gateway-backup
query-positive family and compare against acquisition disabled; they do not establish an optimal
query policy, an adaptivity advantage over a strong local trigger, or universal evidence
acquisition. Third, the WirelessOpsAgent-style arm is a method-contract adaptation over the same
candidate interface, not an official source-code reproduction. Fourth, the held-out Qili/NASA
POWER experiment is source/task transfer in simulation rather than a field deployment. Finally,
the hazard process remains an external Operational Task authority; the paper does not claim a
landslide predictor.

The paper deliberately stops where the frozen workloads stop. A proposed low-cost complete-basis
optimizer was not pursued because 235 formal model requests contained no real alternative-proof
choice. Likewise, task-sensitive validity horizons are already adjacent to AoI/VoI and would
require a workload in which commitment validity, rather than ordinary evidence freshness, changes
the action. We therefore freeze the current method instead of manufacturing additional structure
for algorithmic appearance.

\section{Conclusion}
Intermittent communication turns Agent context management into an execution problem: evidence has
an owner and path, unresolved dependencies become live or retire with candidate plans, and a
correct decision may remain physically in flight after the model call ends. The decision-semantic
compiler makes those relations explicit. On query-negative tasks it preserves the deterministic
physical reference while reducing model cost against a strong same-interface assurance
adaptation; on a query-positive family it acquires real owner evidence, changes plan feasibility,
commits a fallback, and changes the communication outcome with two model families. The resulting
claim is intentionally narrower than generic active learning or semantic communication, but it is
a complete closed loop from Task and Evidence to semantic commitment and asynchronous physical
execution.

\bibliographystyle{IEEEtran}
\bibliography{../../refs}

\end{document}
